% Generated draft: requires mathematical and source review.
\documentclass[11pt]{article}
\usepackage[margin=25mm]{geometry}
\usepackage{fontspec}
\setmainfont{Latin Modern Roman}
\usepackage{amsmath,amssymb,mathtools,unicode-math}
\setmathfont{Latin Modern Math}
\usepackage{enumitem,needspace,xurl,hyperref}
\setlength{\parindent}{0pt}
\setlength{\parskip}{5pt}
\setlength{\emergencystretch}{3em}
\newcommand{\N}{\mathbb{N}}
\newcommand{\Z}{\mathbb{Z}}
\newcommand{\Q}{\mathbb{Q}}
\newcommand{\R}{\mathbb{R}}
\newcommand{\C}{\mathbb{C}}
\newcommand{\F}{\mathbb{F}}
\newcommand{\A}{\mathbb{A}}
\newcommand{\PP}{\mathbb{P}}
\newcommand{\E}{\mathbb{E}}
\newcommand{\eps}{\varepsilon}
\newcommand{\dd}{\,\mathrm{d}}
\begin{document}
\section*{UnsolvedMath ID: 30002972}
\textbf{Source problem code:} catalogue entry
\subsection{Definitions and mathematical statement}
Let \(S\) be a generic projective K3 surface satisfying
\[
\operatorname{Pic}(S)=\mathbb ZH,\qquad H^2=20.
\]
Here \(H^2\) denotes the self-intersection number. Choosing the ample sign of the generator gives a primitive polarization of genus \(1+H^2/2=11\). Write
\[
X=S^{[2]}:=\operatorname{Hilb}^{2}(S)
\]
for the Hilbert scheme parametrizing zero-dimensional closed subschemes of \(S\) of length two, including nonreduced subschemes.

The source does not explicitly specify the ground field. Under the usual complex interpretation, a K3 surface is a smooth connected projective surface over \(\C\) with trivial canonical bundle and \(H^1(S,\mathcal O_S)=0\). In that interpretation, a nontrivial nonsymplectic involution on \(X\) is an automorphism \(\iota\colon X\to X\) satisfying
\[
\iota^2=\operatorname{id}_X,\qquad \iota\ne\operatorname{id}_X,\qquad
\iota^*\sigma=-\sigma
\quad\text{for every }\sigma\in H^0(X,\Omega_X^2).
\]
The complex interpretation is explicit here to explain the terminology; it is not an independently verified hypothesis of the source.

\textbf{Requested statement.} Give a geometric construction or description of the nontrivial nonsymplectic involution on \(S^{[2]}\) referred to in the supplied genus-\(11\) problem, for generic \(S\) with the displayed Picard group and intersection number. Such a description should identify the resulting map with the involution intended by the source.

\textbf{Unresolved qualifications.} The supplied record does not specify the meaning of ``generic,'' identify the intended involution by additional data, or give a formal criterion for a ``geometric description.'' Its definite article presupposes an intended involution; this draft does not assert a separate existence or uniqueness theorem. The ground field and the precise genericity quantifier require confirmation from the report.
\subsection{Short English statement}
For a generic projective K3 surface with Picard group generated by a class of square \(20\), describe geometrically the nonsymplectic involution intended by the source on its Hilbert scheme of length-two subschemes. The ground field, precise meaning of genericity, and identification of the intended involution remain qualifications requiring review. The supplied partially solved label is dated curation, not an independently verified mathematical conclusion.
\subsection{Sources}
\begin{itemize}
\item Supplied local UnsolvedMath registry record, permanent ID \(30002972\), source problem code \texttt{catalogue entry}, in the Oberwolfach Reports Open Problems collection. No individual proposer is supplied; the recorded proposed year is \(2015\).
\item Supplied report citation: \emph{Mini-Workshop: Singular Curves on \(K3\) Surfaces and Hyperkähler Manifolds} (2016), \emph{Oberwolfach Reports}, pp.\ 2939--2967, DOI \texttt{10.4171/owr/2015/51}. \url{https://doi.org/10.4171/owr/2015/51}. The record describes Oberwolfach Reports as documenting mathematical workshops and their open-problem discussions. The cited year \(2016\) and recorded proposed year \(2015\) are retained as supplied.
\item Supplied verification annotation: the statement is labelled \texttt{corrected\_verified}, with a claim that it was checked against the official report and preserves its hypotheses, notation, and target. This annotation is supplied provenance and has not been independently verified here.
\item Supplied literature assessment, checked on \texttt{2026-08-21}: status \texttt{partially\_solved}. The supplied assessment says that the involution is arithmetically characterized and that several geometric constructions, including EPW and related realizations, exist, while no uniform geometric description covers all cases posed in the 2015 question. This is supplied dated curation, not an independently verified assessment; the source itself requests independent verification of time-sensitive status and cited claims.
\item Cattaneo, \emph{Automorphisms of Hilbert Schemes}, as labelled in the supplied literature sources. \url{https://doi.org/10.1002/mana.201800557}.
\item \emph{Geometric Description of Polarized Hilbert Squares}, as labelled in the supplied literature sources. \url{https://doi.org/10.1016/j.jalgebra.2023.01.013}.
\item Quedo, \emph{Geometric Realizations of Non-Symplectic Involutions}, as labelled in the supplied literature sources. \url{https://doi.org/10.1007/s12215-025-01342-1}.
\end{itemize}
\end{document}
